
\section{\BookRecordIndex{T0003}$J$-全纯曲线的奇点}\label{sec:E-7}

现在完成 \BookRecordIndex{T0036}Micallef--White 定理的证明。该定理断言，$J$-全纯奇点不会比代数情形的奇点更坏：穿过给定点的一条 $J$-曲线的各局部分支，可同时通过一个 $C^1$ 共轭变为单复变量多项式的芽。

即使 $M=\mathbb C^n$，且曲线全纯并恰有一个分支，这也并不明显。事实上，设 $u$ 具有形式
\[
  u(z)=(z^k,f(z)),
\]
其中 $f(z)=a_{k+1}z^{k+1}+\cdots$ 是解析函数（可有无穷多个非零系数）。我们的任务是找一个局部微分同胚 $\Psi:(\mathbb C^n,0)\to(\mathbb C^n,0)$，使得
\[
  \Psi\circ u(z)=(z^k,a_{k+1}z^{k+1}+\cdots+a_mz^m),
\]
其中 $m>k$ 是充分大的整数。下文将看到，若 $u$ 为单射，只要 $z\mapsto(z^k,a_{k+1}z^{k+1}+\cdots+a_mz^m)$ 也为单射，就能做到这一点。我们将寻找如下形式的映射：
\[
  \Psi(x,y)=(x,y-\eta(x,y)f_m(x^{1/k})),
\]
其中 $\eta:\mathbb C^n\to[0,1]$ 是在 $u$ 的像上等于 $1$ 的\BookRecordIndex{T0188}截断函数，而
\begin{equation}\label{eq:E-7-1}
  f_m(z)=a_{m+1}z^{m+1}+a_{m+2}z^{m+2}+\cdots.
\end{equation}
除非说明如何选取本来多值的 $x^{1/k}$，否则此映射 $\Psi$ 尚未良定义。

虽然 $\Psi$ 必须延拓为全局 $C^1$ 微分同胚，其定义条件却只涉及 $\operatorname{im}u$ 上的点。这个集合可视为多值函数 $x\mapsto f(x^{1/k})$ 的图像。下文将说明，该函数的不同分支相隔足够远，因而可在各分支附近用不同公式定义 $\Psi$。为更精确地说明这一点，须先作准备。本节始终固定整数 $n\ge2$。对 $\ell\in\mathbb Z$，用 $\mathcal O_\ell$ 表示由偶对 $(U,f)$ 组成的空间，其中 $U\subset\mathbb C$ 是零点的开邻域，而 $f:U\setminus\{0\}\to\mathbb C^{n-1}$ 是满足 (E.5.1) 的光滑映射（见 E.5 节）。我们把 $\mathcal O_\ell$ 的元素视为映射而非芽，简写为 $f\in\mathcal O_\ell$。所以对 $f,g\in\mathcal O_\ell$，记号 $f=g$ 意味着二者定义域相同，且在定义域每一点都相等，而非仅在零点某个邻域上相等。对每个正整数 $k$ 及每个整数 $\ell$，$k$ 次单位根群
\[
  G(k):=\{\zeta\in\mathbb C\mid\zeta^k=1\}
\]
通过下式作用于 $\mathcal O_\ell$：
\[
  G(k)\times\mathcal O_\ell\to\mathcal O_\ell:
  (\zeta,f)\longmapsto f\circ\zeta,
\]
其中 $f\circ\zeta(z):=f(\zeta z)$。

\begin{definition}\label{def:E-7-1}
设 $k$ 为正整数，$\ell,m$ 为整数。若子集 $S\subset\mathcal O_\ell$ 在 $G(k)$ 的作用下不变，则称之为 \emph{\BookRecordIndex{T0005}$k$-不变}。若 $S\subset\mathcal O_\ell$ 为 \BookRecordIndex{T0005}$k$-不变集，且映射 $\Gamma:S\to\mathcal O_m$ 对每个 $\zeta\in G(k)$ 都满足
\[
  \Gamma(f\circ\zeta)=\Gamma(f)\circ\zeta,
\]
则称 $\Gamma$ 为 \emph{\BookRecordIndex{T0006}$k$-等变}。若子集 $S\subset\mathcal O_\ell$ 对任意两个不同元素 $f,g\in S$ 均满足：存在常数 $\varepsilon>0$，使得
\[
  |f(z)-g(z)|\ge\varepsilon|z|^m
\]
对充分小的 $z$ 成立，则称 $S$ 为 \emph{\BookRecordIndex{T0007}$m$-分离}。
\end{definition}

下一个引理处理 $N$ 条曲线 $u_i:=(z^k,f_i(z))$ 的情形。集合 $S_m$ 应被看作相应首阶项集合的 $G(k)$-轨道。

\begin{lemma}\label{lem:E-7-2}
设 $m\ge k\ge1$ 为整数，且 $f_1,\ldots,f_N\in\mathcal O_{k+1}$。设 $S_m\subset\mathcal O_{k+1}$ 是 \BookRecordIndex{T0005}$k$-不变且 \BookRecordIndex{T0007}$m$-分离的集合，$\widetilde f_1,\ldots,\widetilde f_N\in S_m$，满足
\[
  \widetilde f_j-f_j\in\mathcal O_{m+1},\qquad j=1,\ldots,N.
\]
设 $\Gamma:S_m\to\mathcal O_{m+1}$ 为 \BookRecordIndex{T0006}$k$-等变映射，且 $\eta>0$。则存在局部 $C^1$ 映射
\[
  E_\Gamma:(\mathbb C^n,0)\to(\mathbb C^{n-1},0),
\]
满足 $dE_\Gamma(0)=0$，并且当 $x\in\mathbb C$、$y\in\mathbb C^{n-1}$ 充分小时，
\[
  |y|>\eta|x|\quad\Longrightarrow\quad E_\Gamma(x,y)=0,
\]
且对 $j=1,\ldots,N$ 及充分小的 $z$ 有
\begin{equation}\label{eq:E-7-2}
  E_\Gamma(z^k,f_j(z))=\Gamma(\widetilde f_j)(z).
\end{equation}
\end{lemma}

有了这个引理，便可继续寻找所需的微分同胚 $\Psi$。我们只须把引理~\ref{lem:E-7-2} 用于 $N=1$ 和一个映射 $f\in\mathcal O_{k+1}$。若 $f$ 全纯，可取
\[
  S_m:=\{(f-f_m)\circ\zeta^j\mid1\le j\le k\},
  \qquad \Gamma(f-f_m):=f_m,
\]
其中 $f_m$ 表示 $f$ 的高阶项（见式~\eqref{eq:E-7-1}）。这里假设整数 $m$ 取得足够大，使得对每个 $\zeta\in G(k)$，$f\circ\zeta\ne f$ 都蕴含 $(f-f_m)\circ\zeta\ne f-f_m$。设 $E_\Gamma:(\mathbb C^n,0)\to(\mathbb C^{n-1},0)$ 是引理~\ref{lem:E-7-2} 给出的 $C^1$ 映射，定义
\[
  \Psi(x,y):=(x,y-E_\Gamma(x,y)).
\]
此映射在零点附近是局部 $C^1$ 微分同胚，并由~\eqref{eq:E-7-2} 对 $\widetilde f=f-f_m$ 得
\[
  \Psi(z^k,f(z))=(z^k,f(z)-E_\Gamma(z^k,f(z)))
  =(z^k,f(z)-f_m(z)).
\]
故 $\Psi$ 正是去掉 $u$ 的高阶项的微分同胚。

一般情形下 $f$ 不全纯，未必存在这样的截断复多项式 $f-f_m$。这就引出了 Micallef 与 White 证明中的第二个主要思路。对每个 $m>k$，他们在 $f$ 的轨道中取适当映射的平均，构造元素 $A_m^f\in\mathcal O_{k+1}$，使集合
\[
  S_m:=\{A_m^g\mid g\in G(k)\cdot f\}
\]
为 \BookRecordIndex{T0005}$k$-不变且 \BookRecordIndex{T0007}$m$-分离，并且 $f-A_m^f\in\mathcal O_{m+1}$。这便是下一个引理的内容。注意，即使在全纯情形，该引理也能给出更精细的 $u$ 的正规形，即 Puiseux 展开；见下文习题~\ref{ex:E-7-4}，以及 Eisenbud 与 Neumann 的书~\srcite{96} 第 1 章附录。

\begin{lemma}\label{lem:E-7-3}
设 $k$ 为正整数，$S\subset\mathcal O_{k+1}$ 是由定义域相同的映射组成的有限 \BookRecordIndex{T0005}$k$-不变集。假设对每对 $f,g\in S$ 且 $f\ne g$，存在整数 $m(f,g)>k$ 与非零向量 $a_{f,g}\in\mathbb C^{n-1}$，使得
\[
  f(z)-g(z)=z^{m(f,g)}a_{f,g}+\mathcal O_{m(f,g)+1}.
\]
置 $m(f,f):=\infty$ 及 $a_{f,f}:=0$。对 $f\in S$ 与 $m\ge k$，令 $A_m^f$ 表示 $S$ 中与 $f$ 一致到 $m+1$ 阶的映射之平均：
\[
  A_m^f:=\frac1{|S_m^f|}\sum_{g\in S_m^f}g,
  \qquad S_m^f:=\{g\in S\mid g-f\in\mathcal O_{m+1}\}.
\]
则集合 $S_m:=\{A_m^f\mid f\in S\}$ 为 \BookRecordIndex{T0005}$k$-不变且 \BookRecordIndex{T0007}$m$-分离，映射 $S\to S_m:f\mapsto A_m^f$ 为 \BookRecordIndex{T0006}$k$-等变。此外，对一切 $f,g\in S$、$f\ne g$ 及每个整数 $m>k$，有：
\begin{enumerate}
\item[\textup{(i)}] $f-A_m^f\in\mathcal O_{m+1}$。
\item[\textup{(ii)}] 若 $m\ge\max\{m(f,g)\mid g\in S\setminus\{f\}\}$，\sourcefootnote{原书此处的最大值在 $S=\{f\}$ 时取自空集，未给约定；在该情形下，依定义直接有 $S_m^f=\{f\}$ 与 $A_m^f=f$。}则 $A_m^f=f$。
\item[\textup{(iii)}] 若 $m\ge m(f,g)$，则 $A_m^f(z)-A_m^g(z)=z^{m(f,g)}a_{f,g}+\mathcal O_{m(f,g)+1}$。
\item[\textup{(iv)}] $A_m^f=A_m^g$ 当且仅当 $m<m(f,g)$。
\item[\textup{(v)}] 存在向量 $a_m^f\in\mathbb C^{n-1}$，使 $A_m^f(z)=A_{m-1}^f(z)+z^ma_m^f+\mathcal O_{m+1}$。
\end{enumerate}
\end{lemma}

有了此引理，便可在 $f$ 不全纯时如下进行。设 $f\in\mathcal O_{k+1}$，并假设集合 $S:=G(k)\cdot f$ 满足引理~\ref{lem:E-7-3} 的假设。对 $m>k$，定义 $\Gamma_m:S_m\to\mathcal O_{m+1}$：
\[
  \Gamma_m(A_m^f)(z):=A_m^f(z)-A_{m-1}^f(z)-a_m^fz^m,
\]
而 $\Gamma_k(A_k^f):=A_k^f$。由引理~\ref{lem:E-7-3}，$\Gamma_m$ 对所有 $m\ge k$ 均满足引理~\ref{lem:E-7-2} 的假设。关键是当 $m$ 充分大（例如 $m\ge N$）时 $A_m^f=f$ 且 $a_m^f=0$。定义
\[
  \Psi(x,y):=\left(x,y-\sum_{m=k}^{N}E_{\Gamma_m}(x,y)\right).
\]
下文将更详细地解释，公式
\[
  f(z)=A_k^f(z)+\sum_{m=k+1}^{N}
  \bigl(A_m^f(z)-A_{m-1}^f(z)-a_m^fz^m\bigr)
  +\sum_{m=k+1}^{N}a_m^fz^m
\]
蕴含
\[
  \Psi(z^k,f(z))=\left(z^k,\sum_{m=k+1}^{N}a_m^fz^m\right).
\]
这里每个 $E_{\Gamma_m}$ 都起到去掉 $f$ 中对应的高阶项 $A_m^f(z)-A_{m-1}^f(z)-a_m^fz^m$ 的作用。现在证明这两个引理，并说明它们如何\BookRecordIndex{T0326}推出一般形式的定理~\ref{thm:E-1-1}。

\begin{exercise}\label{ex:E-7-4}
设 $f:\mathbb C\to\mathbb C^2$ 为函数 $f(z)=(z^4,z^6+z^8+z^9e^z)$。利用上述过程寻找 $\Psi$，使 $\Psi\circ f(z)=(z^4,z^6+z^9)$。
\end{exercise}

\begin{proof}[引理~\ref{lem:E-7-2} 的证明]
固定常数 $\eta>0$。选取光滑\BookRecordIndex{T0188}截断函数 $\chi:[0,\infty)\to[0,1]$，使
\[
  \chi(r)=\begin{cases}
    1,&r\le1/4,\\
    0,&r\ge1/3.
  \end{cases}
\]
由于 $S_m$ 是 $\mathcal O_{k+1}$ 的有限子集，\sourcefootnote{原书在此使用 $S_m$ 有限，但引理 E.7.2 的陈述只假设 $S_m$ 为 \BookRecordIndex{T0005}$k$-不变、\BookRecordIndex{T0007}$m$-分离集。有限性或足以取得共同定义域、统一分离常数的一致性条件尚需补入；后文具体应用中的 $S_m$ 确实有限。}必要时缩小定义域，可假设每个 $f\in S_m$ 都定义在同一个以零点为中心的开圆盘 $U\subset\mathbb C$ 上。此外，对每个 $\alpha\in\mathbb N^2$，存在常数 $C_\alpha>0$，使
\[
  |\partial^\alpha f(z)|\le C_\alpha|z|^{k+1-|\alpha|}
\]
对每个 $f\in S_m$ 与每个 $z\in U$ 成立。\sourcefootnote{原书把上述估计的幂次写为 $m+1-|\alpha|$，但引理只假设 $S_m\subset\mathcal O_{k+1}$，不能\BookRecordIndex{T0326}推出 $m+1$ 阶的估计（取 $k=1$、$m>1$、$S_m=\{z^2\}$ 即是反例）。此处及下方 $\delta$ 的条件、式 (E.7.5) 的相应幂次据假设改为 $k+1$ 阶；结论不变。}回想 $S_m$ 为 \BookRecordIndex{T0007}$m$-分离集（见定义~\ref{def:E-7-1}）。因此，再缩小 $U$ 后，可假设存在常数 $\varepsilon>0$，使每对 $f,g\in S_m$ 且 $f\ne g$ 均满足
\begin{equation}\label{eq:E-7-3}
  |f(z)-g(z)|\ge\varepsilon|z|^m
\end{equation}
对 $z\in U$ 成立。可将 $\varepsilon$ 取得小于 $\eta$。

控制扰动 $E_\Gamma(x,y)$ 的支集非常重要。为此，注意到由于 $m\ge k$ 且 $\varepsilon<\eta$，可取常数 $\delta>0$，使
\[
  C_0\delta^{1/k}
    +\frac\varepsilon2\delta^{(m-k)/k}<\eta.
\]
这个常数的意义如下：若 $f\in S_m$、$z\in U$、$x\in\mathbb C$ 与 $y\in\mathbb C^{n-1}$ 满足
\begin{equation}\label{eq:E-7-4}
  z^k=x,\qquad |y-f(z)|\le\frac\varepsilon2|z|^m,
  \qquad |x|<\delta,
\end{equation}
则 $|y|\le\eta|x|$；事实上
\begin{equation}\label{eq:E-7-5}
  |y|\le|f(z)|+\frac\varepsilon2|z|^m
  \le\left(C_0|x|^{1/k}
    +\frac\varepsilon2|x|^{(m-k)/k}\right)|x|
  \le\eta|x|.
\end{equation}

给定 $(x,y)\in\mathbb C\times\mathbb C^{n-1}$，$0<|x|<\delta$，若存在满足~\eqref{eq:E-7-4} 的偶对 $(f,z)\in S_m\times U$，则定义
\[
  E_\Gamma(x,y):=
  \chi\!\left(\frac{|y-f(z)|}{\varepsilon|z|^m}\right)
  \Gamma(f)(z).
\]
若不存在这样的偶对（或 $x=0$ 且 $y=0$），则置 $E_\Gamma(x,y):=0$。由~\eqref{eq:E-7-3} 可知，对每个满足 $z^k=x$\sourcefootnote{原书此处写 $z^2=x$；本引理固定任意 $k\ge1$，且条件 (E.7.4) 与下文各处均为 $z^k=x$，故补正。}、$|x|<\delta$ 的 $z\in U$ 及每个 $y\in\mathbb C^{n-1}$，满足~\eqref{eq:E-7-4} 的 $f\in S_m$ 至多有一个。又偶对 $(f,z)$ 满足~\eqref{eq:E-7-4}，当且仅当 $(f\circ\zeta^{-1},\zeta z)$ 也满足。因而，满足~\eqref{eq:E-7-4} 的偶对在商空间 $S_m\times_{G(k)}U$ 中的等价类只取决于 $(x,y)$。由此 $f(z)$ 与 $\Gamma(f)(z)$ 不依赖 $f,z$ 的选择，所以 $E_\Gamma$ 良定义。

当 $|y|\ge\eta|x|$ 且 $|x|<\delta$ 时，$E_\Gamma(x,y)$ 显然为零：此时由~\eqref{eq:E-7-5} 可知，不存在满足~\eqref{eq:E-7-4} 的偶对。由构造也显然可得 $E_\Gamma$ 满足~\eqref{eq:E-7-2}：由于 $\widetilde f_j-f_j\in\mathcal O_{m+1}$，对充分小的 $z\in U$ 有
\[
  |f_j(z)-\widetilde f_j(z)|\le\frac\varepsilon4|z|^m,
\]
故偶对 $(\widetilde f_j,z)$ 在 $x:=z^k$、$y:=f_j(z)$ 时满足~\eqref{eq:E-7-4}，从而
\[
  E_\Gamma(z^k,f_j(z))=\Gamma(\widetilde f_j)(z).
\]

余下须证明 $E_\Gamma$ 连续可微且 $dE_\Gamma(0)=0$。定义开集 $\Omega\subset\mathbb C\times\mathbb C^{n-1}$：\sourcefootnote{原书以下最小值表达式在 $x=0$ 时含分母 $|z|^m=0$，未说明约定；$x=0,y\ne0$ 时可由附近 $x\ne0$ 的零值邻域延拓理解，零点则另行处理。}
\[
  \Omega:=\left\{(x,y)\ \middle|\ |x|<\delta,\ |y|>\eta|x|
  \text{ 或 }\min_{f\in S_m,\ z^k=x}
    \frac{|y-f(z)|}{\varepsilon|z|^m}>\frac13\right\}.
\]
由于 $\chi(r)=0$ 对 $r\ge1/3$ 成立，$E_\Gamma$ 在 $\Omega$ 上为零。因此只须考察 $\Omega$ 补集中的点附近的 $E_\Gamma$。设 $(x_0,y_0)\in\mathbb C\times\mathbb C^{n-1}$，\sourcefootnote{原书此处把 $y_0$ 所在空间印作 $\mathbb C^n$，但本引理的映射定义域及 $y$ 的维数均为 $\mathbb C^{n-1}$，据此补正。}且 $z_0\in U$ 满足
\[
  0<|x_0|<\delta,\quad |y_0|\le\eta|x_0|,
  \quad |y_0-f(z_0)|\le\frac\varepsilon3|z_0|^m,
  \quad z_0^k=x_0.
\]
则在 $x_0$ 附近存在光滑分支 $x\mapsto x^{1/k}$，其在 $x=x_0$ 时取值 $z_0$，并在 $(x_0,y_0)$ 的某邻域内满足
\[
  |y-f(x^{1/k})|<\frac\varepsilon2|x|^{m/k}.
\]
这表明 $E_\Gamma$ 在局部可写为
\begin{equation}\label{eq:E-7-6}
  E_\Gamma(x,y)=
  \chi\!\left(\frac{|y-f(x^{1/k})|}{\varepsilon|x|^{m/k}}\right)
  \Gamma(f)(x^{1/k}),
\end{equation}
因此在原点以外光滑。$E_\Gamma$ 在原点可微且 $dE_\Gamma(0)=0$，由以下估计\BookRecordIndex{T0326}推出：
\begin{equation}\label{eq:E-7-7}
  |E_\Gamma(x,y)|\le c|x|^{(m+1)/k}
\end{equation}
对充分小的 $x$ 成立。常数 $c>0$ 取使
\[
  |\Gamma(f)(z)|\le c|z|^{m+1}
\]
对每个 $f\in S_m$ 及充分小的 $z$ 成立（这是因为 $\Gamma(S_m)\subset\mathcal O_{m+1}$）。式~\eqref{eq:E-7-7} 随即由 $E_\Gamma(x,y)=0$，或 $E_\Gamma(x,y)=\chi(r)\Gamma(f)(z)$ 且 $r\ge0$、$f\in S_m$、$z\in U$、$z^k=x$ \BookRecordIndex{T0326}推出。

再证 $dE_\Gamma$ 在原点连续。这由下列估计\BookRecordIndex{T0326}推出：
\begin{equation}\label{eq:E-7-8}
  |dE_\Gamma(x,y)|\le C\left(|x|^{(m+1-k)/k}+|x|^{1/k}\right)
\end{equation}
对充分小的 $x$ 与适当常数 $C>0$ 成立。为证~\eqref{eq:E-7-8}，在 $x_0\ne0$ 附近选取局部 $k$ 次根 $x\mapsto x^{1/k}$；其导数具有形式 $x\mapsto x^{(1-k)/k}/k$。对~\eqref{eq:E-7-6} 右端求导。项 $\Gamma(f)'(x^{1/k})x^{(1-k)/k}$ 可由 $|x|^{(m+1-k)/k}$ 估计。映射 $y\mapsto |x|^{-m/k}|y-f(x^{1/k})|$ 的导数由 $|x|^{-m/k}$ 控制。在 $\chi$ 的导数不为零处，$|y-f(x^{1/k})|$ 由常数乘以 $|x|^{m/k}$ 控制，故映射 $x\mapsto |x|^{-m/k}|y-f(x^{1/k})|$ 的导数由 $|x|^{-1}+|x|^{(1-m)/k}$ 控制。\sourcefootnote{原书这里只给出 $|x|^{-1}$ 一项；这依赖于前述未被假设保证的 $m+1$ 阶估计。改用 $f\in\mathcal O_{k+1}$ 后还需计入第二项，乘上 $\Gamma(f)$ 的阶数后仍得到式 (E.7.8)。}这些项乘以 $|\Gamma(f)(x^{1/k})|\le c|x|^{(m+1)/k}$ 后，分别给出 $|x|^{1/k}$、$|x|^{(m+1-k)/k}$ 与 $|x|^{2/k}$ 的上界；最后一项可由 $|x|^{1/k}$ 控制。这证明~\eqref{eq:E-7-8}，引理得证。
\end{proof}

\begin{proof}[引理~\ref{lem:E-7-3} 的证明]
先证 (i)：$A_m^f$ 是有限多个满足 $f-g\in\mathcal O_{m+1}$ 的函数 $g\in\mathcal O_{k+1}$ 的平均，因此 $f-A_m^f\in\mathcal O_{m+1}$。

证 (ii)。假设对所有 $g\in S\setminus\{f\}$ 都有 $m\ge m(f,g)$。则 $S_m^f=\{f\}$。事实上，若 $g\in S$ 且 $g\ne f$，依假设 $g-f\notin\mathcal O_{m(f,g)+1}$，从而 $g-f\notin\mathcal O_{m+1}$，即 $g\notin S_m^f$。故 $A_m^f=f$。

证 (iii)。若 $m\ge m(f,g)$，由 (i) 得
\[
  A_m^f-A_m^g-(f-g)\in\mathcal O_{m+1}
    \subset\mathcal O_{m(f,g)+1}.
\]
又依假设 $f-g-z^{m(f,g)}a_{f,g}\in\mathcal O_{m(f,g)+1}$，故
\[
  A_m^f-A_m^g-z^{m(f,g)}a_{f,g}
    \in\mathcal O_{m(f,g)+1},
\]
如所称。

证 (iv)。若 $m<m(f,g)$，则 $f-g\in\mathcal O_{m(f,g)+1}\subset\mathcal O_{m+1}$，所以 $g\in S_m^f$。因而 $S_m^f=S_m^g$，进而 $A_m^f=A_m^g$。若 $m\ge m(f,g)$ 且 $f\ne g$，则由 (iii) 得 $A_m^f-A_m^g\notin\mathcal O_{m(f,g)+1}$，故 $A_m^f\ne A_m^g$。

证 (v)。先注意，对每个 $g\in S$，
\begin{equation}\label{eq:E-7-9}
  g\in S_{m-1}^f\setminus S_m^f
    \quad\Longleftrightarrow\quad m(f,g)=m.
\end{equation}
当 $g=f$ 时两侧显然均不成立。当 $g\ne f$ 时，依假设 $f-g\in\mathcal O_{m(f,g)}\setminus\mathcal O_{m(f,g)+1}$，故 $m=m(f,g)$ 当且仅当 $f-g\in\mathcal O_m\setminus\mathcal O_{m+1}$，亦即 $g\in S_{m-1}^f\setminus S_m^f$。定义
\[
  a_m^f:=\frac1{|S_{m-1}^f|}
     \sum_{m(f,g)=m}a_{f,g},
\]
求和遍历所有满足 $m(f,g)=m$ 的 $g\in S$。由 (i)，可见
\[
  A_m^f-A_{m-1}^f=f-A_{m-1}^f+\mathcal O_{m+1}.
\]
因此
\begin{align*}
  A_m^f-A_{m-1}^f
  &=\frac1{|S_{m-1}^f|}
      \sum_{g\in S_{m-1}^f}(f-g)+\mathcal O_{m+1}\\
  &=z^ma_m^f+\mathcal O_{m+1}.
\end{align*}
最后一个等式是因为由~\eqref{eq:E-7-9}，每个 $g\in S_{m-1}^f\setminus S_m^f$ 都满足 $m(f,g)=m$，而每个 $g\in S_m^f$ 都满足 $m<m(f,g)$。\sourcefootnote{原书此处将后一个不等号印为 $m>m(f,g)$；由 $S_m^f=\{g\mid g-f\in\mathcal O_{m+1}\}$ 及 $m(f,f)=\infty$，应为 $m<m(f,g)$。这也正是相应差在 $m$ 阶消失的条件。}

至此已证 (i)--(v)。再证集合 $S_m$ 为 \BookRecordIndex{T0007}$m$-分离集。若 $f,g\in S$ 且 $A_m^f\ne A_m^g$，则由 (iv) 得 $m\ge m(f,g)$。于是由 (iii)，对充分小的 $z$ 存在常数 $\varepsilon>0$，使
\[
  |A_m^f(z)-A_m^g(z)|
  \ge\varepsilon|z|^{m(f,g)}\ge\varepsilon|z|^m.
\]
映射 $S\to\mathcal O_{k+1}:f\mapsto A_m^f$ 由定义显然为 \BookRecordIndex{T0006}$k$-等变；因此 $S_m$ 为 \BookRecordIndex{T0005}$k$-不变集，引理得证。
\end{proof}

\begin{proof}[定理~\ref{thm:E-1-1} 的证明]
设 $u:\Sigma\to M$ 为一条 \BookRecordIndex{T0003}$J$-全纯曲线，$\Omega\subset\Sigma$ 为开集，满足 $\overline\Omega\subset\Sigma$ 紧，$x_0\in u(\Omega)\setminus u(\partial\Omega)$，并且
\[
  \{z_1,\ldots,z_N\}=u^{-1}(x_0)\cap\Omega.
\]
不失一般性，设 $M=\mathbb C^n$、$x_0=0$，且 $J$ 是 $\mathbb C^n$ 上在原点等于标准复结构的\BookRecordIndex{T0088}殆复结构。对 $j=1,\ldots,N$，令
\[
  T_j=\mathbb Cc_j\subset\mathbb C^n
\]
为 $u_j$ 在 $z=0$ 处的切空间，并设 $u_j$ 在 $z=0$ 的阶为 $k_j\ge1$（见注~\ref{rem:E-5-11}）。于是 $du_j(0)=0$ 当且仅当 $k_j\ge2$。记 $\Pi_j:=T_j^\perp$ 为 $T_j$ 的正交补。由命题~\ref{prop:E-5-10}，存在两两不交的邻域 $U_j\subset\Omega$ 和 $C^2$ 坐标卡 $\varphi_j:(U_j,z_j)\to(\mathbb D,0)$，使 $\varphi_j-\operatorname{id}\in\mathcal O_2$，当 $k_j=1$ 时 $\varphi_j$ 光滑，并且复合映射
\[
  \widetilde u_j:=u\circ\varphi_j^{-1}
\]
具有形式\sourcefootnote{原书此处写 $c_{\nu(j)}$，但本段只引入 $c_j$ 与阶数 $k_j$，未说明 $\nu(j)$ 的含义。这里保留原记号；首项系数须落在 $T_j$ 中，如何同 $c_j$ 精确对应，原文未说明。}
\[
  \widetilde u_j(z)=z^{k_j}c_{\nu(j)}+f_j(z),
\]
其中 $f_j:\mathbb D\to\Pi_j$ 属于 $\mathcal O_{k_j+1}$。先考虑所有切空间相等、所有曲线在零点具有同一阶数的情形。

\noindent\textbf{步骤 1。}固定正整数 $k$。假设每个 $j$ 均满足
\[
  T_j=\mathbb C\times\{0\},\qquad k_j=k.
\]
必要时在 $\mathbb C^n$ 中再作线性坐标变换，并在靶空间中对 $\varphi_j$ 作重标度，可假设
\[
  \widetilde u_j(z)=(z^k,f_j(z)),
\]
其中 $f_j:\mathbb D\to\mathbb C^{n-1}$ 属于 $\mathcal O_{k+1}$。定义 $S\subset\mathcal O_{k+1}$：
\[
  S:=\{f_j\circ\zeta^\ell\mid1\le j\le N,\ 1\le\ell\le k\},
\]
其中 $\zeta(z):=e^{2\pi i/k}z$。注意映射 $f_j\circ\zeta^\ell$ 与 $u_j$ 由同一公式联系，只须把 $\varphi_j$ 换成 $\zeta^{-\ell}\circ\varphi_j$。所以由命题~\ref{prop:E-6-1} 可知，存在一个以零点为中心的开圆盘 $U\subset\mathbb D$，使得：若 $f,g\in S$ 且 $f-g$ 在 $U$ 上不恒为零，则存在整数 $m(f,g)>k$ 和向量 $a_{f,g}\in\mathbb C^{n-1}$，满足
\[
  f(z)-g(z)-z^{m(f,g)}a_{f,g}\in\mathcal O_{m(f,g)+1}.
\]
以下把 $f\in S$ 记为其在 $U$ 上的限制。于是集合 $S$ 满足引理~\ref{lem:E-7-3} 的假设。

对每个 $f\in S$ 与每个整数 $m\ge k$，按引理~\ref{lem:E-7-3} 定义 $A_m^f\in\mathcal O_{k+1}$；当 $m>k$ 时再定义 $a_m^f\in\mathbb C^{n-1}$。\sourcefootnote{原书此处对所有 $m\ge k$ 同时称 $a_m^f$ 已由引理 E.7.3 定义，但该引理 (v) 只在 $m>k$ 时给出此系数；下文 $m=k$ 的定义只用 $A_k^f$。}记
\[
  S_m:=\{A_m^f\mid f\in S\},
\]
并定义映射 $\Gamma_m:S_m\to\mathcal O_{m+1}$：$\Gamma_k(A_k^f):=A_k^f$，以及
\[
  \Gamma_m(A_m^f)(z):=A_m^f(z)-A_{m-1}^f(z)-a_m^fz^m,
  \qquad m>k.
\]
我们证明这个映射良定义。$m=k$ 时由定义即知；当 $m>k$ 且 $A_m^f=A_m^g$ 时，有 $m<m(f,g)$，由引理~\ref{lem:E-7-3}(iv) 得 $A_{m-1}^f=A_{m-1}^g$。再由引理~\ref{lem:E-7-3}(v)，多项式 $z\mapsto z^m(a_m^f-a_m^g)$ 属于 $\mathcal O_{m+1}$，故 $a_m^f=a_m^g$。所以 $\Gamma_m$ 良定义。引理~\ref{lem:E-7-3} 还断言 $S_m$ 为 \BookRecordIndex{T0005}$k$-不变且 \BookRecordIndex{T0007}$m$-分离集。$\Gamma_m$ 的 \BookRecordIndex{T0006}$k$-等变性由其良定义性及映射 $f\mapsto A_m^f$ 的 \BookRecordIndex{T0006}$k$-等变性\BookRecordIndex{T0326}推出。因此，映射 $\Gamma_m:S_m\to\mathcal O_{m+1}$ 对 $\widetilde f_j:=A_m^{f_j}$ 满足引理~\ref{lem:E-7-2} 的全部要求。

对 $m\ge k$，令 $E_{\Gamma_m}:(\mathbb C^n,0)\to(\mathbb C^{n-1},0)$ 为引理~\ref{lem:E-7-2} 断言的局部 $C^1$ 映射。置 $m_0$\sourcefootnote{原书只对 $f\ne g$ 的集合取最大值；当 $S$ 只有一个元素时该集合为空，$m_0$ 未定义。并入 $\{k\}$ 可涵盖此端点，且不改变 $|S|\ge2$ 时的取值。} 为
\[
  m_0:=\max\bigl(\{k\}\cup
    \{m(f,g)\mid f,g\in S,\ f\ne g\}\bigr)
\]
并定义局部 $C^1$ 微分同胚 $\Psi:(\mathbb C^n,0)\to(\mathbb C^n,0)$：
\[
  \Psi(x,y):=\left(x,y-\sum_{m=k}^{m_0}E_{\Gamma_m}(x,y)\right),
\]
其中 $x\in\mathbb C$、$y\in\mathbb C^{n-1}$ 充分小。断言对每个 $f\in S$ 都有
\begin{equation}\label{eq:E-7-10}
  \Psi(z^k,f(z))=
  \left(z^k,\sum_{m=k+1}^{m_0}a_m^fz^m\right).
\end{equation}
事实上，由引理~\ref{lem:E-7-2}，当 $m>k$ 时，
\[
  E_{\Gamma_k}(z^k,f(z))=A_k^f(z),\qquad
  E_{\Gamma_m}(z^k,f(z))
    =A_m^f(z)-A_{m-1}^f(z)-a_m^fz^m.
\]
又由引理~\ref{lem:E-7-3}，
\[
  f(z)=A_k^f(z)+\sum_{m=k+1}^{m_0}
    \bigl(A_m^f(z)-A_{m-1}^f(z)-a_m^fz^m\bigr)
    +\sum_{m=k+1}^{m_0}a_m^fz^m.
\]
这里用到了对每个 $f\in S$ 均有 $A_{m_0}^f=f$。把最后两个等式代入 $\Psi$ 的定义，便得\sourcefootnote{原书下一式第一个求和的下限印为 $m=k+1$，漏掉 $E_{\Gamma_k}=A_k^f$；依 $\Psi$ 的定义及紧邻上式，须从 $m=k$ 起求和。}
\[
  \Psi(z^k,f(z))
  =\left(z^k,f(z)-\sum_{m=k}^{m_0}E_{\Gamma_m}(z^k,f(z))\right)
  =\left(z^k,\sum_{m=k+1}^{m_0}a_m^fz^m\right),
\]
如所称。由于 $f_j\in S$（$j=1,\ldots,N$），由~\eqref{eq:E-7-10} 知 $\Psi$ 与 $\varphi_j$ 在步骤 1 的假设下满足定理~\ref{thm:E-1-1} 的要求。

\noindent\textbf{步骤 2。}假设对每个 $j$ 均有 $T_j=\mathbb C\times\{0\}$。令 $k$ 为各 $k_j$ 的最小公倍数。对函数
\[
  v_j(z):=\widetilde u_j(z^{k/k_j})
\]
应用步骤 1，得到 $C^1$ 微分同胚 $\Psi:\mathbb C^n\to\mathbb C^n$，使得
\[
  \Psi\circ v_j(z)
    =\Psi(z^k,f_j(z^{k/k_j}))=(z^k,P_j(z)),
  \qquad j=1,\ldots,N,
\]
其中 $P_j$ 为如下形式的复多项式：
\[
  P_j(z)=a_{j,k+1}z^{k+1}+\cdots+a_{j,m_0}z^{m_0}.
\]
由 $v_j$ 的定义，$P_j\circ\zeta^{k_j}=P_j$。这表明当 $\ell\notin(k/k_j)\mathbb Z$ 时 $a_{j,\ell}=0$。因此
\[
  P_j(z)=Q_j(z^{k/k_j}),
\]
其中 $Q_j$ 是复多项式。于是对 $j=1,\ldots,N$，
\[
  \Psi(z^{k_j},f_j(z))=(z^{k_j},Q_j(z)).
\]

\noindent\textbf{步骤 3。}考虑一般情形。引理~\ref{lem:E-7-2} 表明，$\Psi$ 的支集（即不被 $\Psi$ 固定的点集之闭包）可选得任意接近 $\mathbb C\times\{0\}$。更精确地，
\[
  \operatorname{supp}(\Psi)
  \subset\{(x,y)\in\mathbb C\times\mathbb C^{n-1}\mid
             |y|\le\eta|x|\},
\]
其中 $\eta>0$ 可任意小。选取有限个两两线性无关的向量 $a_1,\ldots,a_r$，使每个子空间 $T_j$ 都等于某个 $\rho\in\{1,\ldots,r\}$ 对应的 $\mathbb Ca_\rho$。再取常数 $\eta>0$，使去心锥
\[
  V_\rho:=\{\lambda a_\rho+v\mid
       \lambda\in\mathbb C\setminus\{0\},\ v\perp a_\rho,
       \ 0\le|v|\le\eta|\lambda a_\rho|\},
  \qquad\rho=1,\ldots,r,
\]
两两不交。对每个 $\rho$，取酉矩阵 $\Lambda_\rho:\mathbb C^n\to\mathbb C^n$，使 $\Lambda_\rho(a_\rho\mathbb C)=\mathbb C\oplus\{0\}$。

由步骤 2，存在一族 $C^1$ 微分同胚 $\Psi_\rho$，其支集包含在 $\Lambda_\rho(\overline V_\rho)$ 中，且当 $c_j\in\mathbb Ca_\rho$ 时
\[
  \Psi_\rho\circ\Lambda_\rho\circ u_j\circ\varphi_j^{-1}(z)
    =(z^{k_j},z^{k_j}p_j(z)),
\]
其中每个 $p_j$ 都是零点处为零的复多项式。现在置
\[
  \Psi:=\prod_\rho
     \Lambda_\rho^{-1}\circ\Psi_\rho\circ\Lambda_\rho,
\]
并定义
\[
  L_j:=\Lambda_\rho,\qquad c_j\in\mathbb Ca_\rho.
\]
由于去心锥 $V_\rho$ 两两不交，这些映射满足定理~\ref{thm:E-1-1} 的全部条件。
\end{proof}
